\chapter{Transformations de types de donn\'ees}
\label{ch:type-transforms}

\begin{quote}
\emph{<< Les algorithmes de machine learning ne comprennent pas les cat\'egories --- ils comprennent les nombres. >>}
\end{quote}

% --- hook ---
Le machine learning, malgré l'aura mystique qui l'entoure, fait des additions et des multiplications. C'est tout. Or rien dans <<~Marié~>>, <<~Divorcé~>>, <<~Veuf~>>, <<~Célibataire~>> ne se prête à une addition. Encoder, ce n'est donc pas une étape technique préliminaire --- c'est la décision qui transforme une question humaine (état civil) en quelque chose qu'un algorithme peut manipuler. Le choix entre one-hot et target encoding sur une colonne à 10~000 catégories peut faire bouger la performance de plus que le choix d'algorithme lui-même. Ce chapitre traite les transformations de type pour ce qu'elles sont : des décisions de modélisation, pas du nettoyage.
\section{Pourquoi transformer les types ?}

La plupart des algorithmes ML demandent une entr\'ee num\'erique. Les datasets bruts contiennent des variables cat\'egorielles (<< Male >>, << Female >>), des variables ordinales (<< Low >>, << Medium >>, << High >>), et des variables num\'eriques sur des \'echelles tr\`es diff\'erentes (\^age en ann\'ees, revenu en milliers). Un encodage et une mise \`a l'\'echelle propres sont essentiels.

\begin{datatip}
La transformation de type a deux faces : l'\textbf{encodage} convertit les cat\'egories en nombres, la \textbf{mise \`a l'\'echelle} place les nombres sur des plages comparables. Les deux sont n\'ecessaires pour la plupart des pipelines ML.
\end{datatip}

% ============================================================
\section{Encodage cat\'egoriel}

\subsection{Encodage ordinal (label)}

\begin{minted}{python}
import pandas as pd
from sklearn.preprocessing import OrdinalEncoder

# Adult Census
url = ("https://archive.ics.uci.edu/ml/machine-learning-databases/"
       "adult/adult.data")
cols = ["age", "workclass", "fnlwgt", "education", "education_num",
        "marital_status", "occupation", "relationship", "race",
        "sex", "capital_gain", "capital_loss", "hours_per_week",
        "native_country", "income"]
df = pd.read_csv(url, header=None, names=cols, na_values=" ?",
                 skipinitialspace=True)

# Encodage ordinal : pour les variables avec un ordre naturel
edu_order = [["Preschool", "1st-4th", "5th-6th", "7th-8th", "9th",
              "10th", "11th", "12th", "HS-grad", "Some-college",
              "Assoc-voc", "Assoc-acdm", "Bachelors", "Masters",
              "Prof-school", "Doctorate"]]

enc = OrdinalEncoder(categories=edu_order,
                     handle_unknown="use_encoded_value",
                     unknown_value=-1)
df["education_ord"] = enc.fit_transform(df[["education"]])
print(df[["education", "education_ord"]].drop_duplicates()
        .sort_values("education_ord"))
\end{minted}

\begin{warningbox}
L'encodage ordinal impose un ordre num\'erique. L'utiliser sur des cat\'egories nominales (noms de pays par exemple) implique << France $<$ Allemagne $<$ Japon >>, ce qui n'a aucun sens et peut tromper les mod\`eles.
\end{warningbox}

\subsection{One-hot encoding}

\begin{minted}{python}
from sklearn.preprocessing import OneHotEncoder

# One-hot encoding : pour les cat\'egories nominales (sans ordre)
ohe = OneHotEncoder(sparse_output=False, drop="first",
                    handle_unknown="ignore")
encoded = ohe.fit_transform(df[["sex", "race"]])
encoded_df = pd.DataFrame(encoded,
                          columns=ohe.get_feature_names_out())
print(encoded_df.head())
print(f"New columns: {encoded_df.shape[1]}")
\end{minted}

\begin{minted}{python}
# pandas get_dummies : one-hot rapide
df_encoded = pd.get_dummies(df, columns=["sex", "race"],
                            drop_first=True, dtype=int)
print(df_encoded.columns.tolist()[-10:])
\end{minted}

\begin{preprocesstip}
Utilisez \texttt{drop="first"} (ou \texttt{drop\_first=True}) pour \'eviter la multicolin\'earit\'e. Avec $k$ cat\'egories, il ne faut que $k-1$ variables indicatrices --- la cat\'egorie supprim\'ee devient le niveau de r\'ef\'erence.
\end{preprocesstip}

\subsection{Target encoding}

\begin{minted}{python}
# Target encoding : remplacer chaque cat\'egorie par la moyenne de la cible
# Utile pour les colonnes \`a forte cardinalit\'e (100+ cat\'egories)
target_means = df.groupby("occupation")["income"].apply(
    lambda x: (x == ">50K").mean()
).to_dict()

df["occupation_target_enc"] = df["occupation"].map(target_means)
print(df[["occupation", "occupation_target_enc"]]
        .drop_duplicates().sort_values("occupation_target_enc"))
\end{minted}

\begin{warningbox}
Le target encoding cause une fuite de donn\'ees s'il est appliqu\'e avant la s\'eparation train/test. Ajustez toujours sur l'entra\^inement seul, puis transformez les deux. Utilisez \texttt{cross\_val\_predict} ou un target encoding par K-fold pour r\'eduire le surapprentissage.
\end{warningbox}

\subsection{Encodages fr\'equence et binaire}

\begin{minted}{python}
# Encodage par fr\'equence : remplacer la cat\'egorie par sa fr\'equence
freq = df["native_country"].value_counts(normalize=True)
df["country_freq"] = df["native_country"].map(freq)
print(df[["native_country", "country_freq"]].head(10))

# Binary encoding : moins de colonnes que one-hot pour forte cardinalit\'e
# La biblioth\`eque category-encoders fournit cela
# pip install category-encoders
# import category_encoders as ce
# encoder = ce.BinaryEncoder(cols=["native_country"])
# df_binary = encoder.fit_transform(df)
\end{minted}

% ============================================================
\section{Mise \`a l'\'echelle num\'erique}

\subsection{StandardScaler (normalisation z-score)}

\begin{minted}{python}
from sklearn.preprocessing import StandardScaler

scaler = StandardScaler()
df[["age_scaled", "hours_scaled"]] = scaler.fit_transform(
    df[["age", "hours_per_week"]]
)

print(df[["age", "age_scaled", "hours_per_week", "hours_scaled"]].describe())
\end{minted}

\subsection{MinMaxScaler}

\begin{minted}{python}
from sklearn.preprocessing import MinMaxScaler

minmax = MinMaxScaler(feature_range=(0, 1))
df[["age_minmax"]] = minmax.fit_transform(df[["age"]])

# Apr\`es mise \`a l'\'echelle : min=0, max=1
print(df["age_minmax"].describe())
\end{minted}

\subsection{RobustScaler}

\begin{minted}{python}
from sklearn.preprocessing import RobustScaler

# Utilise m\'ediane et IQR : robuste aux outliers
robust = RobustScaler()
df[["capital_gain_robust"]] = robust.fit_transform(df[["capital_gain"]])
print(df["capital_gain_robust"].describe())
\end{minted}

\begin{preprocesstip}
Choisissez votre scaler selon les donn\'ees et le mod\`ele :
\begin{itemize}
  \item \textbf{StandardScaler} : choix par d\'efaut pour la plupart des algorithmes (SVM, r\'egression logistique, r\'eseaux de neurones).
  \item \textbf{MinMaxScaler} : quand on veut des valeurs born\'ees (sigmoid en sortie de r\'eseau).
  \item \textbf{RobustScaler} : quand les donn\'ees contiennent des outliers significatifs.
\end{itemize}
\end{preprocesstip}

% ============================================================
\section{Comparer les scalers visuellement}

\begin{minted}{python}
import matplotlib.pyplot as plt
import numpy as np

fig, axes = plt.subplots(1, 4, figsize=(16, 4))

# Original
df["capital_gain"].hist(bins=50, ax=axes[0])
axes[0].set_title("Original")

# Standard
axes[1].hist(StandardScaler().fit_transform(df[["capital_gain"]]),
             bins=50)
axes[1].set_title("StandardScaler")

# MinMax
axes[2].hist(MinMaxScaler().fit_transform(df[["capital_gain"]]),
             bins=50)
axes[2].set_title("MinMaxScaler")

# Robust
axes[3].hist(RobustScaler().fit_transform(df[["capital_gain"]]),
             bins=50)
axes[3].set_title("RobustScaler")

plt.tight_layout()
plt.savefig("scaler_comparison.png", dpi=150)
plt.show()
\end{minted}

% ============================================================
\section{Discr\'etisation (binning)}

\begin{minted}{python}
# Binning par largeur \'egale
df["age_bin_width"] = pd.cut(df["age"], bins=5,
                              labels=["very_young", "young", "mid",
                                      "senior", "elderly"])

# Binning par fr\'equence (quantiles)
df["age_bin_quantile"] = pd.qcut(df["age"], q=5,
                                  labels=["Q1", "Q2", "Q3", "Q4", "Q5"])

# Bins personnalis\'es selon connaissance m\'etier
df["age_group"] = pd.cut(df["age"],
                          bins=[0, 25, 35, 50, 65, 100],
                          labels=["<25", "25-34", "35-49", "50-64", "65+"])

print(df["age_group"].value_counts().sort_index())
\end{minted}

\begin{minted}{python}
# KBinsDiscretizer de sklearn
from sklearn.preprocessing import KBinsDiscretizer

kbd = KBinsDiscretizer(n_bins=5, encode="ordinal", strategy="quantile")
df[["hours_binned"]] = kbd.fit_transform(df[["hours_per_week"]])
print(df[["hours_per_week", "hours_binned"]].head(10))
\end{minted}

% ============================================================
\section{Transformations de puissance}

\begin{minted}{python}
from sklearn.preprocessing import PowerTransformer

# Yeo-Johnson : marche avec valeurs positives et n\'egatives
pt = PowerTransformer(method="yeo-johnson")
df[["capital_gain_yj"]] = pt.fit_transform(df[["capital_gain"]])

# Box-Cox : ne marche que pour valeurs strictement positives
# pt_bc = PowerTransformer(method="box-cox")

# Transformation log : alternative manuelle pour donn\'ees asym\'etriques
df["capital_gain_log"] = np.log1p(df["capital_gain"])

print(f"Original skew:    {df['capital_gain'].skew():.2f}")
print(f"Yeo-Johnson skew: {df['capital_gain_yj'].skew():.2f}")
print(f"Log skew:         {df['capital_gain_log'].skew():.2f}")
\end{minted}

% ============================================================
\section{Exercices}

\begin{exercise}
\begin{enumerate}
  \item Chargez Titanic. Appliquez un one-hot encoding sur \texttt{Sex} et \texttt{Embarked}. Appliquez un encodage ordinal sur \texttt{Pclass} (1\textsuperscript{re} $>$ 2\textsuperscript{e} $>$ 3\textsuperscript{e}). Montrez le DataFrame r\'esultant.
  \item Pour Adult Census, comparez l'effet de StandardScaler, MinMaxScaler et RobustScaler sur \texttt{capital\_gain}. Tracez les trois distributions.
  \item Impl\'ementez un target encoding sur \texttt{native\_country} avec validation crois\'ee 5-fold pour pr\'evenir la fuite. Comparez aux valeurs d'un target encoding na\"if.
  \item Cr\'eez des groupes d'\^age pour Titanic avec des bins m\'etier (enfant, adolescent, jeune adulte, age moyen, senior). Calculez le taux de survie par groupe d'\^age.
  \item Appliquez une transformation de puissance Yeo-Johnson sur toutes les colonnes num\'eriques d'Adult. Comparez la skewness avant et apr\`es.
\end{enumerate}
\end{exercise}

% ============================================================
\section{R\'esum\'e du chapitre}

\begin{itemize}
  \item L'encodage cat\'egoriel convertit du non-num\'erique en nombres : ordinal pour cat\'egories ordonn\'ees, one-hot pour nominales, target encoding pour forte cardinalit\'e.
  \item La mise \`a l'\'echelle place les variables num\'eriques sur des plages comparables : StandardScaler par d\'efaut, MinMaxScaler pour plages born\'ees, RobustScaler en pr\'esence d'outliers.
  \item La discr\'etisation (binning) convertit les variables continues en cat\'egories, utile pour relations non lin\'eaires et interpr\'etabilit\'e.
  \item Les transformations de puissance (Yeo-Johnson, Box-Cox, log) r\'eduisent l'asym\'etrie et peuvent am\'eliorer la performance des mod\`eles.
  \item Toujours ajuster les transformateurs sur les donn\'ees d'entra\^inement seules, puis les appliquer aux donn\'ees de test.
\end{itemize}
